\subsection{The finite cluster and its critical values}
The quotient used in Lemma~\ref{lem:quot} admits a geometric description. It will matter when one tries to attach an object to an individual factor.

\begin{proposition}\label{prop:cluster}
Fix $t>0$ sufficiently small. Over a neighbourhood of $y^*$, the spectrum of the quantum algebra in Lemma~\ref{lem:quot} is the entire critical scheme in a union of small disjoint neighbourhoods of the critical points of $t\sum_i x^{v_i}$. In this identification the class of $F_t$ corresponds to the quantum Euler element.
\end{proposition}
\begin{proof}
The Fano fan consists of cones over faces of $\Delta=\mathrm{conv}(v_i)$. Every proper face has linearly independent vertices. Thus its face polynomial has no logarithmic critical point, since every monomial coefficient in the vector equation $\sum c_i x^{v_i}v_i=0$ is nonzero. The origin is interior to $\Delta$. The Newton number formula for this convenient nondegenerate-at-infinity potential gives a finite central algebra of length $n!\mathrm{Vol}(\Delta)$ \cite[Lemma~2.8 and proof of Theorem~2.14]{RS}. The simplices joining the origin to the facets have normalized volume one. Their number is the number of fixed points of the dense torus on $X$, hence the Euler characteristic $\mu$. It follows that the central quotient of Lemma~\ref{lem:quot} is an isomorphism.

Choose a small precompact neighbourhood of each central critical point, with no zero of the gradient on its boundary. On their union the relative critical scheme is proper after shrinking the parameter ball, since a critical point cannot cross these compact boundaries. The central fibre is finite, so the relative scheme is finite after a further shrinking. Its $n$ critical equations together with the base coordinates centred at $y^*$ form a system of parameters at every central point of the smooth ambient space. They form a regular sequence. Consequently the finite relative algebra is flat over the base and has rank $\mu$. The invertible coordinate images of Lemma~\ref{lem:quot} map the spectrum of the quantum algebra to this scheme: its finite spectrum lies in the chosen neighbourhoods by properness and the central isomorphism. The induced map between locally free algebras of the same rank is an isomorphism on the central fibre, hence nearby. Choosing the neighbourhoods orbitwise makes this description $G$-equivariant.

For the Euler identification, use the grading on the finite extended mirror before restricting the auxiliary parameters. In the grading with weight zero for the logarithmic volume form and equivariant parameters specialized to zero, the alternative Euler field of \cite[Remark~4.12]{CCIT} is $\sum y_k\partial_{y_k}$. Equivalently, the leading pole of the $z$-direction connection is multiplication by $-F_t$ on the mirror, and by $-[c_1(X)+(1-\mathrm{Gr}_0)\sigma]$ on the quantum side; $\mathrm{Gr}_0$ acts by $j$ on $H^{2j}$. The normalized mirror isomorphism intertwines those poles by \cite[Theorem~4.28 and Remarks~3.2, 4.14]{CCIT}. On the module it takes a monomial times the volume form to $a_k\star c$ at $z=0$, with $c$ the image of the volume form. This is a cyclic generator of the rank-one regular module under the unital algebra comparison, hence a unit. Cancelling it and setting the auxiliary parameters to zero gives
\[
      \sum_{k\in A} y_k a_k=c_1(X)+(1-\mathrm{Gr}_0)\sigma .
\]
This formal identity has convergent coefficients on our finite parameter neighbourhood. The divisor shift in $\tau$ has weight one under $\mathrm{Gr}_0$ and leaves the right-hand side unchanged. Thus a critical value on the selected cluster is the Euler eigenvalue of the corresponding quantum factor.
\end{proof}
